\subsection{Measures}

To balance ethnic diversity with adequate sample sizes, we followed \citeauthor{ma_ethnic_2008}'s
(\citeyear{ma_ethnic_2008}) approach of grouping 55 minority ethnicities into eight larger
categories. This classification reflects population size, geographic concentration,
linguistic and religious affiliations, historical relationships with the central
government, and international sovereignty of ethnic groups:

\begin{enumerate}
    \item Hui: Hui, Salar, Dongxiang, Baoan
    \item Kazakh: Kazakh, Uzbek, Tajik, Kyrgyz, Tatars, Russian
    \item Korean: Korean
    \item Manchu: Manchu, Hezhen, Xibe
    \item Mongolian: Mongolian, Oroqen, Ewenki, Daur
    \item Tibetan: Tibetan, Yugur, Menba (Monba), Lhoba, Tu
    \item Uyghur: Uyghur
    \item Southern minorities: Other ethnic minorities concentrating in Yunnan, Guizhou,
          and Guangxi provinces.
\end{enumerate}

We classified educational attainment into three levels: low (\textit{primary school
    or less}), middle (\textit{middle school}), and high (\textit{high school or above}).
This ensures adequate sample sizes at each educational level across ethnic groups and
census years. For our analysis of status exchange, we created a more granular measure
by converting the original educational scale from 1 (\textit{illiterate}) to 6
(\textit{BA or higher}) into gender-specific percentile ranks within 11-year moving
cohorts.